% !TEX root = ../main.tex

%----------------------------------------------------------------------------------------
% Introduction
%----------------------------------------------------------------------------------------
\setcounter{page}{1} 

\section{Introduction}

\subsection{Learner corpora in Second langauge acquisition}
Collostructional analysis: \citet{stefanowitschCollostructionsInvestigatingInteraction2003}


\begin{quote}
"For studies that aim to contribute to SLA scholarship by providing insights into the progress of L2 acquisition, it will be important to analyze and compare data from learners at different proficiency levels, ideally following a longitudinal design and using corpora that capture data collected from the same learners at different points in time."     
\end{quote}
\citep{learnercorpusRoemer2026}

conditionals more non-native like \citep{Gries_Wulff_2021}


\subsection{Formulaicity in learner language}

going from more formulaic to more productive, increasing variety of verbs that can fill a slot. 
\subsection{Collostructional analysis}


The present study is pseudo-longitudinal in nature and intends to observe, through a collostructional lense, aspects of the development of the second conditional in L2 English learners. 



\begin{questions}
    \item How are type II conditionals dispersed across learner writing tasks and proficiency level in our analyzed data? \label{rq:disp}
    \item Do learners at higher proficiency levels favor different verbs in the protasis of type II conditionals than beginner learners? \label{rq:dca}   
    \item Does the verbal slot in the protasis of type II conditionals provide evidence for a transition from more formulaic to more productive use of the construction? \label{rq:lexical_div}
\end{questions}



